% !TEX root = ../../main.tex
\chapter{Primjeri}
\label{ch:primjeri}

\section{Matematika}

	Okruženje \texttt{align} poravnava redove po znaku \texttt{\&}, koji ide ispred
	znaka jednakosti.
	Kada izraz pređe u novi red, operator se piše na početku novog reda.

	Za konturu $C$ koja ograničava polukružni prsten $1 \le r \le 2$, $0 \le \vp
		\le \pi$ je, po Greenovoj teoremi,
	\begin{align}
		\oint_{C} y^{2} \dd{x} + 3xy \dd{y}
		 & = \iint_{D} \left( \pdv{(3xy)}{x} - \pdv{(y^{2})}{y} \right) \dd{x} \dd{y}
		= \iint_{D} y \dd{x} \dd{y}
		\nonumber                                                                     \\
		 & = \int_{1}^{2} \int_{0}^{\pi} r^{2} \sin\vp \dd{r} \dd{\vp}
		= \frac{r^{3}}{3} \bigg|_{1}^{2} \, (-\cos\vp) \bigg|_{0}^{\pi}
		= \frac{14}{3}.
		\label{eq:green}
	\end{align}
	Jednačina se referencira s \verb+\eqref+, npr.~rezultat~\eqref{eq:green}.
	Ispred svake reference ide \verb+~+, da se broj ne bi odvojio od prethodne
	riječi pri prelomima.

	Paket \texttt{slashed} daje Feynmanovu slash notaciju~\cite{peskin}:
	\begin{equation}
		\Tr\left[\gamma^{\mu}\slashed{p}_{2}\gamma^{\nu}\slashed{p}_{1}\right]
		= 4\left[p_{1}^{\mu}p_{2}^{\nu} + p_{2}^{\mu}p_{1}^{\nu}
			- g^{\mu\nu}(p_{1} \cdot p_{2})\right].
		\label{eq:trag}
	\end{equation}

	Vektorske operacije se pišu kako je uobičajeno kod nas, $\ddiv \rrot \vb{F} =
		0$, a hiperbolni tangens kao $\tgh x$.

\section{Mjerne jedinice}

	Paket \texttt{siunitx} pravilno piše veličine s mjernim jedinicama.
	Pogrešno: $g = 9.81 ms^{-2}$, 12\%.
	Pravilno: $g = \SI{9.81}{\meter\per\second\squared}$, \SI{12}{\percent}.
	Brojevi bez jedinica idu kroz \verb+\num+, npr.~\num{1.23e-4}, a raspon kroz
	\verb+\SIrange+, npr.~\SIrange{88}{94}{\giga\electronvolt}.

\section{Slike}

	Slike se navode bez ekstenzije i putanje ako su u folderu \texttt{slike/}, kao
	na slici~\ref{fig:pmf}.
	Opis slike ide ispod slike, a \verb+\label+ uvijek nakon \verb+\caption+.

	\begin{figure}[h]
		\centering
		\includegraphics[width=0.3\textwidth]{pmf}
		\caption{Logo Prirodno-matematičkog fakulteta.}
		\label{fig:pmf}
	\end{figure}

\section{Feynmanovi dijagrami}

	Paket \texttt{tikz-feynman} crta dijagrame direktno u \LaTeX-u, kao na
	slici~\ref{fig:feynman}.

	\begin{figure}[h]
		\centering
		\begin{tikzpicture}
		\begin{feynman}
			\vertex (a);
			\vertex [above left=of a] (i1) {\(e^{+}\)};
			\vertex [below left=of a] (i2) {\(e^{-}\)};
			\vertex [right=of a] (b);
			\vertex [above right=of b] (f1) {\(f\)};
			\vertex [below right=of b] (f2) {\(\bar{f}\)};

			\diagram* {
				(i2) -- [fermion] (a) -- [fermion] (i1),
				(a) -- [photon, edge label=\(\gamma\)] (b),
				(f2) -- [fermion] (b) -- [fermion] (f1),
			};
		\end{feynman}
	\end{tikzpicture}
		\caption{Anihilacija $e^{+}e^{-} \to f\bar{f}$ preko fotona.}
		\label{fig:feynman}
	\end{figure}

\section{Grafici}

	Stil \texttt{grafik} definisan u \texttt{preamble.tex} daje jednak izgled svim
	graficima.
	Podaci se čitaju iz \texttt{.dat} fajla, a funkcija se može nacrtati direktno,
	kao na slici~\ref{fig:grafik}.

	\begin{figure}[h]
		\centering
		\begin{tikzpicture}
		\begin{axis}[
				grafik,
				xmin=87.5, xmax=94.5, ymin=0,
				xlabel={$E_{\text{CM}}$ /\si{\giga\electronvolt}},
				ylabel={$\sigma$ /\si{\nano\barn}},
			]
			\addplot[blue, domain=87.5:94.5, samples=200]
			{41.5*1.556/((x-91.19)^2+1.556)};
			\addlegendentry{Breit--Wigner}
			\addplot[red, only marks, mark=star, mark size=2.6pt]
			table {chapters/primjeri/podaci.dat};
			\addlegendentry{Podaci}
		\end{axis}
	\end{tikzpicture}
		\caption{Presjek u okolini Z-pola.}
		\label{fig:grafik}
	\end{figure}

	Općenito slobodni ste da ubacujete i iz \texttt{matplotlib} kao obični grafik,
	ali naporno je uštimavati veličinu i stil fonta.

\section{Tabele}

	Tabele koriste \texttt{booktabs}, bez vertikalnih linija.
	Opis tabele ide iznad tabele.
	Kolone tipa \texttt{S} poravnavaju brojeve po decimali, a \texttt{table-format}
	zadaje broj cifara prije i poslije tačke (tabela~\ref{tab:primjer}).

	\begin{table}[h]
		\centering
		\caption{Primjer tabele s brojevima poravnatim po decimalnoj tački.}
		\label{tab:primjer}
		\begin{tabular}{@{}lS[table-format=2.5]S[table-format=2.5]S[table-format=1.4]@{}}
			\toprule
			           & {Proračun} & {Standardni model} & {Greška (\%)} \\
			\midrule
			$R_{\mu}$  & 20.725     & 20.735             & 0.0476        \\
			$R_{\tau}$ & 20.725     & 20.781             & 0.2688        \\
			\addlinespace
			$R_{c}$    & 0.17067    & 0.17221            & 0.8935        \\
			$R_{b}$    & 0.21955    & 0.21581            & 1.7342        \\
			\bottomrule
		\end{tabular}
	\end{table}

\section{Citiranje i izvori}

	Izvori se unose u \texttt{literatura.bib}, a citiraju s \verb+\cite+,
	npr.~knjiga o \TeX-u~\cite{texbook} ili članak~\cite{Bender_1998}.
	Bibliografija prikazuje samo citirane izvore, poredane redom citiranja.

	Svaka Wikipedia stranica~\cite{wiki:latex} ima opciju \texttt{Tools -> Cite
		this page} (slika~\ref{fig:wiki-cite}), na kojoj se nalazi gotov BibTeX kod
	(slika~\ref{fig:wiki-bib}).
	Slično je i s arXivom (slika~\ref{fig:arxiv-cite}).
	U polju \texttt{author} zarez razdvaja prezime od imena, a riječ \texttt{and}
	razdvaja autore.

	\begin{figure}[h]
		\centering
		\begin{subfigure}[t]{0.3\textwidth}
			\centering
			\includegraphics[width=\linewidth]{wikipedia-cite}
			\caption{Cite this page.}
			\label{fig:wiki-cite}
		\end{subfigure}
		%
		\hfill
		\begin{subfigure}[t]{0.65\textwidth}
			\centering
			\includegraphics[width=\linewidth]{wikipedia-bibtex}
			\caption{Kod koji se kopira u \texttt{literatura.bib}.}
			\label{fig:wiki-bib}
		\end{subfigure}
		\caption{Citiranje s Wikipedije.}
	\end{figure}

	\begin{figure}[h]
		\centering
		\includegraphics[width=\textwidth]{arxiv-cite}
		\caption{Citiranje s arXiva.}
		\label{fig:arxiv-cite}
	\end{figure}

	\paragraph{Napomena.}
	Ovo je \verb+\paragraph+, naslov koji se nastavlja u istom redu.
	Koristan je za kratke cjeline koje ne trebaju biti u sadržaju.

\section[Bespotrebno dug naslov]{Ovo je jako dug naslov koji kada bi ovako stajao u sadržaju bilo bi previše, dolazilo bi do preloma, i čitava stvari bi bila u 3 reda}
	Srećom, ako koristimo \verb+\section[Kratki naziv]{Puni naziv}+ to se ne
	dešava.
	Sličan pristup radi i sa slikama i tabelama,
	\verb+\caption[Kratki opis]{Detaljni opis}+, gdje je to puno teže izbjeći.
